\section{Data}
\label{sec:data}

\paragraph{Holdings.} We use the SEC's Form 13F data sets: 54 files, one per
release (quarterly through 2023, then three-month windows). Together they
contain every 13F report filed from July 2013 to August 2026, covering
holdings for 2013Q2 to 2026Q2. Institutional managers with investment
discretion over at least \$100 million in Section 13(f) securities (mainly
U.S.-listed stocks and options) must report their long positions in these
securities within 45 days of each quarter's end. A report covers the whole
management firm, not a single fund. For brevity we call each 13F filer a
``fund''; each N-PORT series is a single fund. We keep common shares only (no
options or bonds), take the latest restatement of each report plus amendments
that add positions, and convert values to dollars (the SEC switched from
thousands to dollars in January 2023).

\paragraph{Samples.} No free list says which 13F filers are hedge funds, so
we build several samples (Table~\ref{tab:samples}).
\emph{Hand-picked hedge funds:} 49 well-known managers across multi-strategy,
quantitative, long/short equity, activist, event-driven and macro styles
(listed in the code repository). We choose them with written rules: we use
the entity that actually reports holdings (related entities often file notices
only), join histories across changes of legal entity, and exclude market
makers, family offices and mostly long-only managers.
\emph{Form ADV hedge funds:} a list chosen in 2026 may favor survivors, so we
also select funds by rule from SEC Form ADV filings made between July 2012 and
June 2013, prior to the start of our sample period. We add up the gross asset value of the
private funds each registered adviser classifies as hedge funds (excluding
feeder funds and funds of funds). We keep advisers whose hedge funds hold at
least \$1 billion, at least half of their private-fund assets and at least a
quarter of their regulatory assets. We link each one to its 13F filings by the
CIK it reports or by exact legal name, require at least 10 holdings in June
2013, and keep the 100 largest, of which 99 hold stocks in our universe. The
list follows each adviser's 2013 filer, so it does not follow later
reorganizations (e.g.\ SAC Capital into Point72).
\emph{Long-only managers:} with the same timing and linking, we select
advisers that draw more than 75\% of their assets from registered investment
companies (mutual funds), have essentially no hedge funds (under 5\% of
regulatory assets), and manage at least \$1 billion. They must also be active
rather than index-like: active share~\cite{cremers2009active} of at least 0.6
against the value-weighted stock universe in June 2013. This excludes, for
example, Vanguard and Fidelity's firm-wide filing, whose combined holdings are
close to the market. Three sponsors of sector and leveraged ETFs also pass.
No ETF flag exists before 2018 (Form N-CEN), so we keep the rule and also
report results without them. Like hedge fund filings, these are firm-level
reports, so the comparison holds the level of aggregation fixed.
\emph{Individual mutual funds:} Cohen et al.~\cite{cohen2026mimicking} study
individual funds, so to compare with them we use the SEC's Form N-PORT data
sets, which report each registered fund's holdings from late 2019. As of
December 2019 we keep active U.S.\ equity funds (at least 20 universe stocks
making up at least 60\% of net assets, active share of at least 0.6, not ETFs)
and take the 300 largest by net assets. To isolate the effect of aggregation,
we also add these funds up to their registrant (the fund trust), for
registrants with at least two selected funds; a family-quarter exists only if
all its funds reported.

\paragraph{Stock universe.} Each quarter we take the 500 largest stocks by
total institutional holdings across all 13F filers. We measure holdings as
each filer's shares times the stock's median reported price per share, and
exclude filers whose own implied price is more than five times off (a unit
error). We also exclude debt securities, ETFs, funds and ADRs (on average
\univtype{} of the 500 largest holdings each quarter) and stocks held by
fewer than 100 filers (another \univholders). Since it is built solely from SEC
data, the universe keeps companies that were later acquired or
delisted; using the S\&P~500 instead would lose 10--28\% of early members to
ticker changes that free data cannot resolve. Membership for quarter $t$ is
fixed at the end of quarter $t-1$.

\paragraph{Market data and news.} We map CUSIPs to tickers with OpenFIGI
(with a share-class fallback for securities listed only abroad). We download
daily dividend- and split-adjusted prices and split histories from Yahoo
Finance for 3,104 tickers, with prices through 30 June 2026 (30 October 2026
for the return tests), downloaded in November 2026. Yahoo's history can
change, so the code repository archives the panel. Company-tagged news
headlines come from FNSPID~\cite{dong2024fnspid} (1.4 million headlines for
676 universe stocks, 2013--2023), scored with FinBERT~\cite{araci2019finbert}
(\texttt{ProsusAI/finbert}). Policy uncertainty is the monthly index
of~\cite{baker2016epu}, and factor returns are the Fama--French five factors
and momentum from Kenneth French's data library.
